% GENERATED FILE. Edit canonical lessons, then run npm run generate:books.
% canonical-source-hash: fnv1a64:52b651f36090212d
% canonical-lessons: ES-C06-trabajar

\chapter{Trabajar --- Work, and the Torture Device Behind It}
\label{ch:spanish-trabajar}
\hlchaptermodality{19}

\begin{chapteropening}
\textbf{By the end of this chapter:} \emph{I can say what I do for work with a second regular -ar verb.}

tripalium, a three-stake device, is why work and travail are the same word.
\end{chapteropening}

\section[trabajar]{trabajar — \textquotedblleft{}to work,\textquotedblright{} and it once meant \emph{torture}}

\label{lesson:ES-C06-trabajar}



\begin{quote}
A second \textbf{-ar} verb — snap it into the exact template you just learned. And \emph{trabajar} has the darkest, best etymology in this chapter: it began as a word for \textbf{torture}, and it's the secret origin of English \textbf{travel}.
\end{quote}

\begin{sounds}
\begin{itemize}
\raggedright
  \item \texttt{j-jota} — the \textbf{j} is a raspy \emph{h}-from-the-throat (the \emph{jota}): \emph{trabajar} = \emph{tra-ba-\textbf{HAR}}.
  \item \texttt{r-tap}, \texttt{b-soft} — tapped \emph{r}, soft \emph{b}.
\end{itemize}
\end{sounds}

\begin{cousinweb}
\textbf{trabajar} comes from Vulgar Latin \textbf{tripaliāre}, \emph{\textquotedblleft{}to torture\textquotedblright{}} — from \textbf{tripalium}, a Roman instrument of torture made of \textbf{three stakes} (\emph{tri-} \textquotedblleft{}three\textquotedblright{} + \emph{pālus} \textquotedblleft{}stake,\textquotedblright{} the root of English \textbf{pale} and \textbf{impale}).

The chain of meaning is grimly logical: \emph{torture} $\to$ \emph{torment / suffering} $\to$ \emph{hard toil} $\to$ just \textbf{work}. And English took the very same word down a parallel road:

\begin{itemize}
\raggedright
  \item \textbf{travail} — painful, laborious effort (and the pains of childbirth).
  \item \textbf{travel} — yes: journeying was once such an \emph{ordeal} that \textquotedblleft{}travail\textquotedblright{} (toil) became \textbf{travel} (to make a journey). Every time you \emph{travel}, you're etymologically being \emph{tortured}.
\end{itemize}

So even plain \textbf{trabajo}, “I work,” hides the old idea of toil as torment — a dark little joke inside an everyday verb.
\end{cousinweb}

\begin{grammarlens}[title={Grammar Lens: same -ar template}]
Drop \textbf{-ar}, add the ending — identical to \emph{hablar}:

\noindent
\begin{tabularx}{\linewidth}{@{}>{\raggedright\arraybackslash}X>{\raggedright\arraybackslash}X>{\raggedright\arraybackslash}X@{}}
\toprule
\textbf{person} & \textbf{form} & \textbf{meaning} \\
\midrule
I (ending carries the subject) & \textbf{trabajo} & I work \\
tú & \textbf{trabajas} & you work \\
usted & \textbf{trabaja} & you (formal) work \\
\bottomrule
\end{tabularx}

Notice you didn't learn a new ending — that's the point. \textbf{One template, every regular -ar verb.}
\end{grammarlens}

\subsection*{Your turn}
Say these aloud:

\begin{itemize}
\raggedright
  \item \textquotedblleft{}trabajar\textquotedblright{} — \emph{tra-ba-HAR}, raspy \emph{j}
  \item \textquotedblleft{}trabajo, trabajas, trabaja\textquotedblright{} — same endings as \emph{hablo/hablas/habla}
  \item \textquotedblleft{}trabajar\textquotedblright{} $\to$ travail $\to$ travel — torture became a journey
\end{itemize}

\begin{tcolorbox}[breakable,colback=teal!4,colframe=teal!35!black,title={Before you move on}]
What did \emph{trabajar} originally mean, and from what device? (\textquotedblleft{}To torture,\textquotedblright{} $\leftarrow$ \emph{tripalium}, a three-stake tool.) What two English words come from the same root? (\emph{Travail} and \emph{travel}.) How do you say \textquotedblleft{}I work\textquotedblright{}? (\emph{Trabajo} — same \emph{-o} ending.) Next: \textbf{estudiar}, \textquotedblleft{}to study.\textquotedblright{}
\end{tcolorbox}
